\documentclass{article}
\usepackage[utf8]{inputenc}
\usepackage{hyperref}
\usepackage{listings}
\usepackage{xcolor}
\usepackage{enumitem}

\lstset{
    basicstyle=\ttfamily\small,
    breaklines=true,
    frame=single,
    backgroundcolor=\color{gray!5},
}

\title{Sort directories by disk usage and display the largest ones}
\author{Marcos de Carvalho}
\date{2026-04-21}

\begin{document}

\maketitle

\begin{abstract}
Displays the top 10 directories with the largest disk usage in the current directory, sorted in descending order by size in human-readable format.
\end{abstract}

\section*{Language}
bash

\section*{Category}
disk-usage

\section*{Command}
\begin{lstlisting}
du -sh */ | sort -hr | head -10
\end{lstlisting}

\section*{Explanation}
The command uses 'du -sh' to calculate the total disk usage for each directory in human-readable format. The '*/' glob pattern matches only directories (not files). The output is piped to 'sort -hr' to sort in descending order using human-readable number comparison, with the largest directories appearing first. Finally, 'head -10' limits output to the top 10 directories.

\section*{Tags}
du, disk-usage, directories, storage, sort, system-administration

\section*{Dependencies}
coreutils


\section*{Arguments}
\begin{enumerate}
  \item \textbf{PATH} (Optional): Directory path to analyze (default: current directory) \\ Default: .
  \item \textbf{COUNT} (Optional): Number of largest directories to display (default: 10) \\ Default: 10
\end{enumerate}

\section*{Examples}
\begin{enumerate}
  \item \texttt{du -sh */ | sort -hr | head -10} - Show top 10 largest subdirectories in current directory
  \item \texttt{du -sh */ | sort -hr} - Show all subdirectories sorted by size without limit
  \item \texttt{du -sh */ | sort -hr | head -5} - Show top 5 largest subdirectories
  \item \texttt{du -sh ./* | sort -hr | head -10} - Include hidden files and directories (those starting with dot)
  \item \texttt{cd /var \&\& du -sh */ | sort -hr | head -10} - Find largest subdirectories in /var
\end{enumerate}

\section*{Output}
\begin{lstlisting}
847G	cache/
512G	data/
256G	backups/
128G	logs/
96G	documents/
64G	downloads/
32G	pictures/
16G	videos/
8.5G	archives/
2.1G	temp/
\end{lstlisting}

\section*{Notes}
\begin{itemize}
  \item The '*/' pattern matches only directories; use '*' to include files as well
  \item The -h flag makes output human-readable (K, M, G, T); use -B for a specific block size
  \item The -s flag summarizes each directory's total size instead of listing contents recursively
  \item Hidden directories (starting with .) are excluded; use './*' or '*' (with shopt) to include them
  \item On systems with many subdirectories, results are calculated quickly since du only traverses each directory once
  \item Symbolic links are not followed by default; add -L flag to follow symlinks
  \item Results may vary if some directories are mounted on different filesystems or are inaccessible due to permissions
\end{itemize}

\section*{Warnings}
\begin{itemize}
  \item Results may be incomplete or misleading if you lack read permissions on some directories
  \item On network filesystems (NFS, SMB), the command may be slower due to network latency
  \item If dealing with millions of directories, the command can consume considerable memory and time
  \item Sparse files may cause inaccurate size reporting on some filesystems
\end{itemize}

\section*{See Also}
\begin{itemize}
  \item find-largest-storage-users
  \item find-large-files-recursive
\end{itemize}

\section*{Status}
reviewed

\section*{Safety}
safe

\section*{Shell}
bash

\section*{Platforms}
linux, gnu-linux, freebsd, openbsd, netbsd

\section*{Created At}
2026-04-21

\section*{Updated At}
2026-04-21

\end{document}
